\section{问题分析与总体思路}

本题的难点不在于直接套用一个时间序列模型，而在于解释“为什么低弹性供需冲击会给出很高理论价格，而附件真实价格却停留在约 110--120 美元/桶的平台附近”。因此，本文将建模思路分成三条互相支撑的证据线。

\begin{enumerate}
  \item 先建立传统供需基准，用题面低短期弹性和供应中断范围计算无缓冲理论上界，作为后续动态模型的反事实参照。
  \item 再建立短期冲击模型，把 SPR、商业库存、绕道运输、需求收缩、恐慌衰减和预期修复纳入同一日度递推框架，解释附件冲突窗口内的真实价格路径。
  \item 最后建立中长期油价调节模型，并通过三情景预测、敏感性分析、蒙特卡洛情景树、滞后地缘风险指数检验和历史样本稳健性检验识别影响长期价格平衡点和二次跳涨风险的关键变量。
\end{enumerate}

\begin{figure}[htbp]
\centering
\begin{tikzpicture}[
  node distance=0.82cm and 0.72cm,
  stage/.style={draw=OilBorder, fill=PaperNoteBack, rounded corners=1pt, align=center, minimum height=0.78cm, text width=2.85cm, font=\small\bfseries},
  note/.style={draw=OilBorder, fill=white, rounded corners=1pt, align=center, minimum height=0.62cm, text width=2.85cm, font=\scriptsize},
  arrow/.style={-{Latex[length=2.2mm]}, line width=0.45pt, draw=OilGray}
]
\node[stage] (data) {附件价格数据};
\node[stage, right=of data] (base) {传统供需基准};
\node[stage, right=of base] (short) {短期冲击模型};
\node[stage, right=of short] (long) {中长期油价调节模型};
\node[note, below=of base] (counter) {低弹性反事实压力};
\node[note, below=of short] (check) {基准对比\\消融与稳健性};
\node[note, below=of long] (risk) {三情景\\蒙特卡洛与状态转移};
\draw[arrow] (data) -- (base);
\draw[arrow] (base) -- (short);
\draw[arrow] (short) -- (long);
\draw[arrow] (base) -- (counter);
\draw[arrow] (short) -- (check);
\draw[arrow] (long) -- (risk);
\draw[arrow] (counter.east) -- (check.west);
\draw[arrow] (check.east) -- (risk.west);
\end{tikzpicture}
\caption{本文总体技术路线}
\label{fig:paper-roadmap}
\end{figure}

这一技术路线的核心原则是：真实价格只来自附件数据；长期部分采用情景预测和区间表达；模型检验覆盖朴素基准、滞后复制、参数过拟合和结构稳定性等关键边界。

最终论文使用赛题给定的两个模型名称组织主线：任务 1 为短期冲击模型，任务 2 为中长期油价调节模型。传统供需模型用于提供反事实基准；三情景、蒙特卡洛和状态转移情景树是中长期油价调节模型内部的条件外推方法；消融实验、敏感性分析和蒙特卡洛情景树用于证明因素是否值得保留，并刻画多因素同时恶化时的尾部风险。因素纳入原则见\cref{tab:factor-inclusion}。

\begin{table}[htbp]
\centering
\caption{因素纳入原则}
\label{tab:factor-inclusion}
\begin{tabularx}{\linewidth}{lX X}
\toprule
层级 & 因素 & 处理方式 \\
\midrule
短期冲击模型物理层 & 供应中断、短期弹性、SPR、商业库存、绕道运输、需求收缩、恐慌衰减 & 进入短期冲击模型，参与 0--60 天路径拟合与机制解释 \\
短期冲击模型价格形成层 & 地缘风险溢价、冲击不确定性、持续封锁制度风险、缓冲确认折价、预期修复折价 & 进入短期冲击模型，但必须通过基准对比、消融实验和稳健性检验证明必要性 \\
中长期调节层 & OPEC+ 剩余产能、非 OPEC 供给响应、冲突升级概率 & 不强行用于短期拟合，进入中长期油价调节模型和尾部风险解释 \\
改进方向 & 油轮保险费、美元指数、期货期限结构、投机资金和套保需求 & 当前缺少附件内真实外生数据，不单独参数化，作为后续扩展方向 \\
\bottomrule
\end{tabularx}
\end{table}

\subsection{可复现实现原则}

本文将工程实现转化为建模可信度的一部分：数据处理保证口径可追溯，仿真程序保证模型可复现，图表系统保证论证过程可检查。具体而言，本文所有清洗数据、外部校验数据、参数校准结果、蒙特卡洛样本、敏感性结果、论文图表和最终 PDF/DOCX 均由统一项目脚本生成，使每一个结论都可以回溯到明确的数据文件、模型参数和图表来源。

\begin{table}[htbp]
\centering
\caption{可复现实现对建模质量的作用}
\label{tab:engineering-implementation}
\begin{tabularx}{\linewidth}{lX X}
\toprule
实现维度 & 在本文中的落地方式 & 对建模质量的作用 \\
\midrule
数据处理 & 自动清洗附件价格序列，接入 EIA、JODI、OPEC、GPR、OVX 等外部校验数据，并统一输出中文命名数据表。 & 减少手工处理误差，保证真实价格来自附件，外部数据只用于数量级校验与长期约束。 \\
仿真与批量计算 & 使用程序完成多目标参数搜索、局部稳健性复核、机制消融、蒙特卡洛情景树和状态转移路径抽样。 & 将“参数是否偶然”“长期是否只是一条线”等问题转化为可重复的数值检验。 \\
图表与论文生成 & 由脚本统一生成拟合图、机制图、敏感性图、蒙特卡洛图、气泡图、热力图和最终论文文件。 & 让图表承担论证功能，而不是只做装饰；每张主图都对应一个关键模型问题。 \\
\bottomrule
\end{tabularx}
\end{table}

\keypoint{本文的实现重点是把数据、仿真和图表组织成可追溯的证据链：每个模型结果都能从原始数据、参数文件和生成脚本中复现。}

\subsection{本文贡献与边界}

围绕赛题目标，本文重点回答三个可检验问题：第一，在题面给定低弹性和封锁冲击下，传统供需模型为什么会明显高估价格；第二，哪些缓冲机制能够把理论高位压回附件中的现实价格平台；第三，当封锁延长至 60--180 天时，哪些因素决定价格中枢和尾部跳涨风险。本文的贡献与边界见\cref{tab:contribution-boundary}。

\begin{table}[htbp]
\centering
\caption{本文贡献与边界}
\label{tab:contribution-boundary}
\begin{tabularx}{\linewidth}{lX X}
\toprule
维度 & 本文贡献 & 识别边界 \\
\midrule
反事实基准 & 依据赛题低弹性假设构造无缓冲价格上界，给出动态模型需要解释的压力来源。 & 该上界是题面条件下的反事实压力，不是对现实市场价格的直接预测。 \\
短期冲击模型 & 将 SPR、商业库存、绕道运输、需求收缩、恐慌衰减、风险溢价和预期修复放入同一日度递推框架，解释附件冲突窗口内的价格平台。 & 适用范围限定于 2026 冲突窗口下的机制解释，不外推为跨年份通用交易模型。 \\
模型鲁棒性检验 & 与朴素上一日、漂移随机游走和 ARIMA 基准比较，并加入 DM 检验、滞后平移、机制消融和参数扰动测试。 & 检验提供支持性证据，但不能替代未来真实观测对长期外推的事后验证。 \\
中长期油价调节模型 & 把短期冲击模型延伸到 60--180 天，输出乐观、中性、悲观路径、蒙特卡洛区间和状态转移尾部风险。 & 中长期结果是条件情景预测，不把中性曲线写成未来每日油价的精确点预测。 \\
外部合理性校验 & 使用 EIA 周度数据、JODI 多国月度石油数据、OPEC 全球供需平衡表、GPR 指数、OVX 隐含波动率和 2017--2026 历史同长度窗口验证数量级、事件极端性与风险环境。 & 外部数据用于数量级校验和约束；JODI 按多国上报口径解释，附件价格仍是短期拟合目标。 \\
\bottomrule
\end{tabularx}
\end{table}

\keypoint{本文的论证链条是“反事实压力—机制缓冲—鲁棒性检验—条件外推”：先说明高价压力来源，再说明缓冲机制如何改变价格路径，并用基准、消融、敏感性和历史窗口检验支撑模型解释。}

\subsection{文献依据与建模约束}

本文将文献调研转化为可执行的建模约束。Kilian（2009）和 Hamilton（2009）表明，油价冲击不能只由当期供应缺口解释，还需要区分供给、需求、预防性需求和宏观环境；Kilian 与 Murphy（2014）进一步说明库存和投机性需求会影响油价路径。因此，本文将传统供需模型定位为低弹性反事实上界，而不是最终预测模型。

在预测评价方面，Alquist 与 Kilian（2010）以及 Alquist、Kilian 与 Vigfusson（2013）指出，原油价格预测必须与无变化预测、随机游走等强基准比较，并明确预测期限和不确定性边界。因此，本文在短期部分设置上一日价格朴素基准、漂移随机游走和滚动 ARIMA；在长期部分不把中性线写成逐日真实价格预测，而是同时给出三情景中心路径、蒙特卡洛扇形区间和状态转移尾部风险。

在政策缓冲方面，Newell 与 Prest（2017）、Kilian 与 Zhou（2019）和 Stevens 与 Zhang（2021）均表明，SPR 释放能够影响油价，但其作用通常是临时削峰和稳定预期，不能被写成完全抵消供应中断的万能机制。因此，本文把 SPR、商业库存和绕道运输都作为部分缓冲变量：它们改变剩余缺口和价格峰值，却不直接决定长期均衡价格。

在风险变量方面，Caldara 与 Iacoviello（2022）提供了 GPR 指数的理论和数据基础，Mei 等（2020）、Liu 等（2019）说明地缘风险更适合解释油价波动和尾部风险。由于 GPR 来自新闻文本，本文将其作为外部风险校验和滞后检验变量，不进入短期同步拟合。进一步地，本文接入 Cboe/FRED 的 OVX 原油 ETF 隐含波动率，将更接近市场交易定价的风险变量用于长期风险权重和不确定性强度约束，而非短期方向预测因子。

\begin{table}[htbp]
\centering
\caption{核心文献对本文模型设计的约束}
\begin{tabularx}{\linewidth}{lX X}
\toprule
文献主线 & 对本文的约束 & 落地位置 \\
\midrule
油价冲击分解 & 不能只用供应缺口解释油价，需纳入库存、预期和风险重估 & 传统基准与短期冲击模型 \\
随机游走强基准 & RMSE 必须和 no-change、Random Walk、ARIMA 等基准比较 & 短期模型基准对比与 DM 检验 \\
SPR 与库存政策缓冲 & SPR 是削峰和争取时间，不是完全填补缺口的万能变量 & 物理供需层、消融实验和长期敏感性分析 \\
GPR 与 OVX 波动风险 & GPR 适合做风险校验，OVX 适合做市场隐含波动率约束，二者均不宜同步解释同日或同月油价方向 & 滞后 GPR 检验、OVX 滞后波动检验和长期尾部风险解释 \\
长期预测不确定性 & 中长期结果应报告情景、区间和尾部概率，不夸大单点预测 & 中长期油价调节模型、蒙特卡洛和状态转移情景树 \\
\bottomrule
\end{tabularx}
\end{table}
